\documentclass[11pt,a4paper]{article}
\usepackage[margin=2.2cm]{geometry}
\usepackage[T1]{fontenc}
\usepackage[utf8]{inputenc}
\usepackage{lmodern}
\usepackage{booktabs}
\usepackage{array}
\usepackage{xcolor}
\usepackage{enumitem}
\usepackage{hyperref}
\usepackage{microtype}
\hypersetup{colorlinks=true,linkcolor=blue!60!black,citecolor=blue!60!black,urlcolor=blue!60!black,breaklinks=true}
\setlist{nosep,leftmargin=*}
\newcommand{\model}{\textsc{TomaNet}}
\newcommand{\modelcls}{\textsc{TomaNet-C}}
\newcommand{\modeldet}{\textsc{TomaNet-D}}

\title{\textbf{Research Context}\\[4pt]
\large Why This Project Exists, and What Current Research Shows}
\author{Tomato disease detection project}
\date{\today}

\begin{document}
\maketitle
\tableofcontents
\newpage

\section{Why This Problem Matters}
Tomato is one of the most widely grown vegetable crops in the world, and diseases like
bacterial spot, blight, and viral infections can wipe out a large share of a farmer's
yield if caught late. A camera and a small model that can spot disease early -- ideally
on an ordinary phone, not a data-centre GPU -- is genuinely useful. That is the practical
goal behind this project: not just a high accuracy number on a lab dataset, but something
that still works on a real photo, in a real field.

\section{What Current Research (2024--2026) Shows}

\subsection{Lab datasets are a solved problem -- and that's the trap}
The most commonly used tomato dataset, PlantVillage, is already essentially solved.
Recent papers report 99.4--99.96\% accuracy on it \cite{salman2025}. The differences
between these papers are smaller than normal run-to-run noise. In plain terms: getting a
slightly higher number on PlantVillage does not prove a new idea is any good, because
everyone is already near the ceiling.

The real problem is that PlantVillage photos are misleadingly easy: they are single
leaves, photographed in a lab, on a plain, uniform background. Researchers have shown
that a classifier looking at \emph{only 8 background pixels} (no leaf at all) can still
guess the right disease 49\% of the time on this dataset, when random guessing should
give about 2.6\% \cite{noyan2022}. The very first big tomato-disease paper, back in 2016,
already found this exact trap: 99.35\% accuracy in the lab, dropping to 31.4\% on photos
from the open web \cite{mohanty2016}. It keeps happening: a 99.96\%-accurate model drops
to 68--74\% the moment it's tested on a different, real-world dataset
\cite{salman2025}. A 2024 review that looked specifically at 35 studies using only real
field photos concluded that models trained on lab data ``perform poorly on real-world
images'' \cite{jelali2024}.

\textbf{What this means for us:} we do not report PlantVillage accuracy as our main
result. We use it as one dataset among several, and we care much more about how a model
trained on it performs on a genuinely different, real field photo.

\subsection{Detection has moved to real field data, and the numbers are public}
For the detection side of the problem, current best practice is a fair comparison of
YOLO-family detectors on the same field dataset, same settings, same seed. On the
Tomato-Village field dataset (the same one we use), the best published result is
mAP@0.5 = 0.936 with a large model, and mAP@0.5 = 0.864 with the smallest, phone-sized
model \cite{ramos2025}. These numbers are the honest target for anyone working at the
same small model size -- not something to casually claim beating without actually
measuring it the same way.

Small tweaks to an existing detector (adding an attention module, changing the loss
function) commonly get reported as +1 to +4 mAP improvements on a single dataset and a
single run \cite{chen2025yolov8n,meng2025}. That size of gain is close to normal
run-to-run noise, so the research community is increasingly sceptical of papers that only
show one such number without repeating the experiment.

\subsection{Nobody has built one small model for both jobs}
A few research groups have built models that do more than one job at once -- for
example, classifying the plant species \emph{and} detecting disease in the same
model. But every example found is heavy: one recent one uses about 49 million
parameters \cite{pmjdm} -- roughly 15--35$\times$ larger than our biggest model. No
lightweight model family (under about 5 million parameters) has been shown to work for
\emph{both} whole-leaf classification and box-level detection on tomato specifically.
This is the most direct gap our project is aimed at.

\subsection{Explaining a model's decision is still just pictures, not numbers}
Many tomato papers show a heat-map (Grad-CAM or similar) next to a photo, as a picture,
to argue ``look, the model is focusing on the right spot.'' What is missing across the
field is a \emph{number} that says how well the heat-map actually lines up with where the
disease really is. This is a known open gap, not something we have solved yet -- but it
shapes what a more rigorous version of this project should measure later, rather than
trusting pictures alone.

\section{What Approach We Are Taking, Because of This}
Given the above, our approach is:
\begin{enumerate}
  \item \textbf{Build one small, shared model for both tasks} (classification and
    detection) -- directly targeting the gap that no lightweight tomato model has done
    this (\S2.3). This is the model described in the companion document, \emph{Our
    Approach}.
  \item \textbf{Never trust a single-dataset, single-run number.} Because the field is
    full of small, single-seed ``improvements'' that don't replicate, we treat any result
    from one dataset or one run as provisional, not a finding.
  \item \textbf{Test on real field data, not just lab photos}, and be explicit about the
    gap between the two, instead of only reporting the easy lab number.
  \item \textbf{Fix data problems before trusting any accuracy number.} The literature
    itself flags leakage and inconsistent splits as a recurring, under-reported problem.
    We found and fixed exactly this kind of issue in our own detection dataset (see the
    companion document for details) before training anything on it.
\end{enumerate}

\section{What We Are Not Claiming}
To stay honest and in scope:
\begin{itemize}
  \item We do not claim state-of-the-art PlantVillage accuracy -- it is saturated and
    not a meaningful contribution on its own.
  \item We do not claim to beat the best published detection numbers on Tomato-Village
    without actually running the same fair comparison ourselves, at the same model size.
  \item Larger research questions raised by the literature -- learning from unlabelled
    images, detecting a disease the model has never seen before, and turning
    explainability into a real number instead of a picture -- are recognised as genuine,
    open gaps, but are future work, not something this project has done yet.
\end{itemize}

\begin{thebibliography}{9}
\footnotesize
\bibitem{mohanty2016} S.~P. Mohanty, D.~P. Hughes, M.~Salath\'e, ``Using deep learning for image-based plant disease detection,'' \emph{Front. Plant Sci.} 7:1419, 2016. doi:10.3389/fpls.2016.01419
\bibitem{noyan2022} M.~A. Noyan, ``Uncovering bias in the PlantVillage dataset,'' arXiv:2206.04374, 2022.
\bibitem{salman2025} Salman, Muhammad, Han, ``Vision transformer with mixture-of-experts for plant disease,'' \emph{Front. Plant Sci.}, 2025. doi:10.3389/fpls.2025.1522985
\bibitem{jelali2024} M.~Jelali, ``Deep learning for tomato disease detection using real-field datasets: a review,'' \emph{Front. Plant Sci.}, 2024. doi:10.3389/fpls.2024.1493322
\bibitem{ramos2025} L.~Ramos, A.~D. Sappa, ``Benchmarking YOLOv8--YOLOv12 for tomato disease detection,'' \emph{Scientific Reports}, 2025. doi:10.1038/s41598-025-11064-0
\bibitem{chen2025yolov8n} M.~Chen et al., ``Improved YOLOv8n for tomato disease detection in the field,'' \emph{Scientific Reports}, 2025. doi:10.1038/s41598-025-00405-8
\bibitem{meng2025} Meng, Chen, Dong, Wang, ``YOLO11n with multi-scale fusion for tomato disease detection,'' \emph{Plants} 14(20):3174, 2025. doi:10.3390/plants14203174
\bibitem{pmjdm} Fu, Wang, Wang, Sun, ``PMJDM: joint species classification and disease detection,'' \emph{Front. Plant Sci.}, 2025. doi:10.3389/fpls.2025.1599671
\end{thebibliography}

\end{document}
